\section*{第 \ref{ch:g10:orders-of-magnitude} 章——数量级：测量宇宙}

\begin{solution}{exo:g10:orders-of-magnitude:1}
\emph{1.} $\qty{750}{nm} = 750 \times \qty{e-9}{m} = \qty{7.50e-7}{m}$.
\emph{2.} $\qty{2.5}{km} = \qty{2.5e3}{m}$.
\emph{3.} $\qty{3.2}{\micro m} = \qty{3.2e-6}{m}$.
\emph{4.} $\qty{0.15}{Tm} = 0.15 \times \qty{e12}{m} = \qty{1.5e11}{m}$
——一个\omterm{def:g10:orders-of-magnitude:lightyear}{天文单位}。
\end{solution}

\begin{solution}{exo:g10:orders-of-magnitude:2}
\emph{1.} $\qty{149600000}{km} = \qty{1.496e11}{m}$：$4$ 位
\omterm{def:g10:orders-of-magnitude:sigfig}{有效数字}。
\emph{2.} $\qty{1.06e-10}{m}$：$3$ 位\omterm{def:g10:orders-of-magnitude:sigfig}{有效数字}（前导零
永不计入）。
\emph{3.} $\num{8.1e9}$：$2$ 位\omterm{def:g10:orders-of-magnitude:sigfig}{有效数字}。
\emph{4.} $\qty{5.20e-3}{s}$：$3$ 位\omterm{def:g10:orders-of-magnitude:sigfig}{有效数字}（末尾的
零是一种“承诺”）。
\end{solution}

\begin{solution}{exo:g10:orders-of-magnitude:3}
\omterm{def:g10:orders-of-magnitude:scinot}{尾数}和\omterm{def:g10:orders-of-magnitude:scinot}{指数}分别计算
（依据\cref{prop:g10:orders-of-magnitude:powers}）：
\[
3 \times 2 \times 10^{8-5} = \num{6e3}, \qquad
\tfrac{8}{4} \times 10^{4-(-3)} = \num{2e7},
\]
\[
5^{2} \times 10^{12} = \num{2.5e13}, \qquad
\tfrac{4 \times 3}{6} \times 10^{-3-9+7} = \num{2e-5} .
\]
\end{solution}

\begin{solution}{exo:g10:orders-of-magnitude:4}
人类：$1.7 < 5$，\omterm{def:g10:orders-of-magnitude:oom}{数量级} $\qty{e0}{m}$。珠穆朗玛峰：
$\qty{8848}{m} = \qty{8.848e3}{m}$，且 $8.848 \geq 5$，\omterm{def:g10:orders-of-magnitude:oom}{数量级}
为 $\qty{e4}{m}$。细菌：\omterm{def:g10:orders-of-magnitude:oom}{数量级} $\qty{e-6}{m}$。地球：
$5.97 \geq 5$，\omterm{def:g10:orders-of-magnitude:oom}{数量级} $\qty{e25}{kg}$。
\end{solution}

\begin{solution}{exo:g10:orders-of-magnitude:5}
\emph{1.} 半个刻度：$L = (127.0 \pm 0.5)\,\unit{mm}$。
\emph{2.} $\Delta L / L = 0.5/127 \approx 0.4\%$.
\emph{3.} 橡皮：$0.5/12 \approx 4\%$。铅笔的测量
精度约高十倍——同一把尺，\omterm{def:g10:orders-of-magnitude:uncertainty}{绝对
不确定度}相同，被测物体长十倍。
\end{solution}

\begin{solution}{exo:g10:orders-of-magnitude:6}
\emph{1.} 比邻星：
$4.24 \times \qty{9.47e15}{m} = \qty{4.01e16}{m}$。天狼星：
$8.6 \times \qty{9.47e15}{m} = \qty{8.1e16}{m}$（两位\omterm{def:g10:orders-of-magnitude:sigfig}{有效
数字}，与输入数据相同）。
\emph{2.} $t = d/c = \qty{2.4e13}{m} / \qty{3.00e8}{m/s}
= \qty{8.0e4}{s}$，即 $\num{8.0e4}/3600 \approx \qty{22}{h}$：
几乎要等一整天才能收到探测器发来的每条消息。
\end{solution}

\begin{solution}{exo:g10:orders-of-magnitude:7}
$t = d/c$，每次都如此。月球：
$\qty{3.84e8}{m}/\qty{3.00e8}{m/s} = \qty{1.28}{s}$。太阳：
$\qty{1.496e11}{m}/\qty{3.00e8}{m/s} = \qty{499}{s} \approx
\qty{8.3}{min}$。海王星：
$\qty{4.50e12}{m}/\qty{3.00e8}{m/s} = \qty{1.50e4}{s} \approx
\qty{4.2}{h}$。你看到海王星时，看到的是它大约
四小时前的样子——每台望远镜都是一台温和的“时光机”。
\end{solution}

\begin{solution}{exo:g10:orders-of-magnitude:8}
\emph{1.} 一张纸：$\qty{52.0}{mm}/500 = \qty{0.104}{mm}
= \qty{1.04e-4}{m}$。不确定度也相应除以 500：
$\qty{0.5}{mm}/500 = \qty{0.001}{mm}$，因此
$t = (104 \pm 1)\,\unit{\micro m}$。
\emph{2.} $1/104 \approx 1\%$.
\emph{3.} 直接测量一张纸，得到
$(0.1 \pm 0.5)\,\unit{mm}$：\omterm{def:g10:orders-of-magnitude:uncertainty}{相对不确定度}约为 $500\%$
——尺子甚至无法证明这张纸存在。将 $500$
张纸叠放后，测得的长度乘以 $500$，而尺子的
\omterm{def:g10:orders-of-magnitude:uncertainty}{绝对不确定度}保持不变，\omterm{def:g10:orders-of-magnitude:uncertainty}{相对不确定度}缩小了
同样的 $500$ 倍。
\end{solution}

\begin{solution}{exo:g10:orders-of-magnitude:9}
将各长度除以 $\qty{2e-10}{m}$。细菌：
$\qty{2.0e-6}{m}/\qty{2e-10}{m} = \num{1.0e4}$ 个原子。头发：
$\qty{7.0e-5}{m}/\qty{2e-10}{m} = \num{3.5e5}$。人体：
$1.7/\num{2e-10} = \num{8.5e9}$，\omterm{def:g10:orders-of-magnitude:oom}{数量级}为 $10^{10}$：从头到脚
约一百亿个原子。
\end{solution}

\begin{solution}{exo:g10:orders-of-magnitude:10}
\emph{1.} $s = \qty{0.22}{m}/\qty{1.39e9}{m} = \num{1.58e-10}$.
\emph{2.} 地球：
$\qty{1.27e7}{m} \times \num{1.58e-10} = \qty{2.0e-3}{m} =
\qty{2.0}{mm}$——一粒胡椒粒——以距足球
$\qty{1.496e11}{m} \times \num{1.58e-10} = \qty{23.7}{m} \approx
\qty{24}{m}$ 的距离绕行。
\emph{3.} 月球：
$\qty{3.84e8}{m} \times \num{1.58e-10} = \qty{6.1}{cm}$，距
胡椒粒。海王星：
$\qty{4.50e12}{m} \times \num{1.58e-10} \approx \qty{710}{m}$：模型
无法放进校园里。
\end{solution}

\begin{solution}{exo:g10:orders-of-magnitude:11}
\emph{1.} $21.0 \times 29.7 = \qty{623.7}{cm^2}$，四舍五入至输入数据的三位
\omterm{def:g10:orders-of-magnitude:sigfig}{有效数字}：$\qty{624}{cm^2}$。
\emph{2.} $v = \qty{100.0}{m}/\qty{12.3}{s} = \qty{8.13}{m/s}$（三位
\omterm{def:g10:orders-of-magnitude:sigfig}{有效数字}，由时间决定）。
\emph{3.} 若取五位\omterm{def:g10:orders-of-magnitude:sigfig}{有效数字}，就等于声称面积精确到
$\qty{0.01}{cm^2}$，而毫米级的输入数据从未
达到过此精度——多写的小数位依然是造假。
\end{solution}

\begin{solution}{exo:g10:orders-of-magnitude:12}
\emph{1.} 半径 $r = \qty{0.55}{m}$，面积
$A = \pi r^{2} = \pi \times 0.55^{2} \approx \qty{0.95}{m^2}$。体积
$V = \qty{1.0}{mm^3} = \qty{1.0e-9}{m^3}$，而薄膜的体积为（面积）
$\times$（厚度），因此
\[
e = \frac{V}{A} = \frac{\qty{1.0e-9}{m^3}}{\qty{0.95}{m^2}}
= \qty{1.1e-9}{m} .
\]
\emph{2.} 分子的尺度约为 $\qty{e-9}{m}$，比
原子的 $\qty{e-10}{m}$ 高一个量级，与分子是小型
原子集群这一点相符。富兰克林的一茶匙油使半英亩池塘
平静了下来，而他自己并未察觉，其实已经测量了纳米尺度。
\end{solution}

\begin{solution}{exo:g10:orders-of-magnitude:13}
假设：静息时每分钟约 $70$ 次心跳，寿命 $80$
年。每天：$70 \times 60 \times 24 \approx \num{1e5}$ 次心跳。每
年：$\num{1e5} \times 365 \approx \num{4e7}$。一生：
$\num{4e7} \times 80 \approx \num{3e9}$ ——\omterm{def:g10:orders-of-magnitude:oom}{数量级}为
$10^{9}$，约几十亿次心跳。若取四位\omterm{def:g10:orders-of-magnitude:sigfig}{有效数字}，
就等于假装清楚八十年里每分钟的心率：
输入值都只是\omterm{def:g10:orders-of-magnitude:oom}{数量级}的估计，输出结果的精度
不可能超过其输入（见\cref{met:g10:orders-of-magnitude:sigfig}）。
\end{solution}

\begin{solution}{exo:g10:orders-of-magnitude:14}
对折 $n$ 次后，厚度为 $2^{n} \times \qty{1.0e-4}{m}$。
\emph{1.} 需 $2^{n} \geq 8848/\num{1.0e-4} = \num{8.8e7}$。由于
$2^{26} \approx \num{6.7e7}$ 不够，而
$2^{27} \approx \num{1.3e8}$ 足够，因此需对折 $27$ 次（纸堆
高度 $\approx \qty{13}{km}$）。
\emph{2.} 需 $2^{n} \geq \num{3.84e12}$：$2^{41} \approx
\num{2.2e12}$ 不够，$2^{42} \approx \num{4.4e12}$ 则足够——
需对折 $42$ 次，纸堆高 $\qty{4.4e8}{m}$，远超月球距离。
\emph{3.} 只能靠实验：每次对折面积减半，
对折七次后，纸堆的厚度便超过其宽度。倍增的算术
本身并无问题——\omterm{def:g10:orders-of-magnitude:scinot}{指数}增长只是超出了纸张
（乃至直觉）所能把握的范围。
\end{solution}

\begin{solution}{exo:g10:orders-of-magnitude:15}
\emph{1.} 月球：$\qty{0.03}{m}/\qty{3.84e8}{m} = \num{7.8e-11}$。
桌子：$\qty{0.001}{m}/\qty{1.000}{m} = \num{1.0e-3}$。月球
测量的精度高出的倍数为
$\num{1.0e-3}/\num{7.8e-11} \approx \num{1.3e7}$ ——超过一千万倍。
\emph{2.} 若在一米长的桌子上达到 $\num{7.8e-11}$ 的精度，意味着
$\Delta x = \qty{7.8e-11}{m}$，小于一个原子的
尺度（$\approx \qty{e-10}{m}$）：桌子的边缘甚至
无法如此\emph{界定}。相对而言，地月
激光测距是人类拥有的最精密的测量之一。
\end{solution}

\begin{solution}{pb:g10:orders-of-magnitude:1}
\textbf{1.} 沙粒：$\qty{0.5}{mm} = \qty{5e-4}{m}$，属于毫米
尺度；原子：$\qty{2.5e-10}{m} = \qty{0.25}{nm}$，属于纳米尺度。

\textbf{2.} $\qty{5e-4}{m}/\qty{2.5e-10}{m} = \num{2.0e6}$：沿直径
排列着两百万个原子。

\textbf{3.} 棱长上各有 $2 \times 10^{6}$ 个原子的立方体，共含
$(2\times10^{6})^{3} = 8 \times 10^{18}$ 个原子，\omterm{def:g10:orders-of-magnitude:oom}{数量级}为
$10^{19}$。

\textbf{4.} $V = (\qty{5e-4}{m})^{3} = \qty{1.25e-10}{m^3}$，因此
$m = \rho V = \qty{2500}{kg/m^3} \times \qty{1.25e-10}{m^3}
= \qty{3.1e-7}{kg} = \qty{0.31}{mg}$。

\textbf{5.} 每个原子的质量：
$\qty{3.1e-7}{kg}/(8\times10^{18}) = \qty{3.9e-26}{kg}$ ——恰好
落在“几 $10^{-26}\,\unit{kg}$”的范围内。简单的原子立方体
再现了真实的原子质量，该模型通过了合理性检验
（\cref{met:g10:orders-of-magnitude:sanity}），且表现优异。

\textbf{6.} $s = \qty{5e-4}{m}/\qty{1.39e9}{m} = \num{3.6e-13}$。

\textbf{7.} 模型地球：直径
$\qty{1.27e7}{m} \times \num{3.6e-13} = \qty{4.6e-6}{m}$ ——相当于
一个细菌的大小，地球距沙粒的距离为
$\qty{1.496e11}{m} \times \num{3.6e-13} = \qty{5.4}{cm}$：
地球是一个距沙粒五厘米的微生物。

\textbf{8.} $\qty{4.5e12}{m} \times \num{3.6e-13} = \qty{1.6}{m}$：整个
行星系统的轨道都在一张餐桌之内。

\textbf{9.} $\qty{4.0e16}{m} \times \num{3.6e-13} = \qty{1.4e4}{m}
= \qty{14}{km}$。

\textbf{10.} 距我们这个桌面版太阳系最近的恒星，是另一粒沙子，
远在十四千米之外。两者之间几乎
空无一物：银河系绝大部分是空旷的空间，这也是为什么
恒星之间的碰撞几乎闻所未闻。

\textbf{11.} $t = \qty{5e-4}{m}/\qty{3.00e8}{m/s} = \qty{1.7e-12}{s}$
——不足两皮秒。

\textbf{12.} 月球：$\qty{3.84e8}{m}/\qty{3.00e8}{m/s} =
\qty{1.28}{s}$。太阳：$\qty{1.496e11}{m}/\qty{3.00e8}{m/s} =
\qty{499}{s} \approx \qty{8.3}{min}$：太阳光
总是八分钟前发出的。

\textbf{13.} 一年是 $365.25 \times 86400 = \qty{3.156e7}{s}$，因此
$\qty{1}{ly} = \qty{3.00e8}{m/s} \times \qty{3.156e7}{s}
= \qty{9.47e15}{m}$；然后
$\qty{4.0e16}{m}/\qty{9.47e15}{m} = 4.2$：数据中的距离确实
约为 $\qty{4.2}{ly}$。

\textbf{14.} $\qty{2.4e22}{m}/\qty{9.47e15}{m} = \num{2.5e6}$：约
两百五十万\omterm{def:g10:orders-of-magnitude:lightyear}{光年}。那道光约在 $2.5$ 百万年前离开仙女座，
那时人属（\emph{Homo}）的第一批成员
正在学习打制石器，而当时世界上
还没有\emph{Homo sapiens}（智人）。

\textbf{15.} $\qty{3.00e8}{m/s} \times \qty{1e-9}{s} = \qty{0.30}{m}$.
在三分之一纳秒内，信号至多能传播约
$\qty{10}{cm}$，因此机器中任意两个必须在一个时钟周期内
相互通信的部件，彼此间距都须在十厘米以内：
速率限制了尺寸。

\textbf{16.} 水中原子的平均质量：
$\qty{3.0e-26}{kg}/3 = \qty{1.0e-26}{kg}$。原子数：
$\qty{70}{kg}/\qty{1.0e-26}{kg} = \num{7e27}$。

\textbf{17.} 以沙粒为标尺：
$\num{7e27}/\num{8e18} \approx 10^{9}$——你所含的原子相当于十亿
粒沙子所含的原子。以恒星为标尺：
$\num{7e27}/\num{2e23} = \num{3.5e4}$——你的原子数超过
可观测宇宙中全部恒星的数目，比值为三万五千比一。

\textbf{18.} $\text{（你）}/\text{（原子）} =
\qty{1.7}{m}/\qty{2.5e-10}{m} \approx \num{7e9} \approx 10^{10}$，
而 $\text{（宇宙）}/\text{（你）} =
\qty{8.8e26}{m}/\qty{1.7}{m} \approx \num{5e26}$。在这一尺度阶梯上，你
大约比原子高十级，比宇宙低近二十七级，
按倍数计，人类非常接近自身的原子。

\textbf{19.} 由条件可得
$x^{2} = 10^{-15} \times 10^{26.9} = 10^{11.9}$，于是
$x = 10^{5.95} \approx \qty{9e5}{m} \approx \qty{900}{km}$——几乎
恰好等于最大的小行星谷神星的直径
（$\approx \qty{940}{km}$）。在原子核与宇宙之间的阶梯正中，
坐落着一颗并不大的矮行星。

\textbf{20.} $\qty{8.8e26}{m} \times \num{3.6e-13} = \qty{3.2e14}{m}$，
即 $\qty{3.2e14}{m}/\qty{1.496e11}{m} \approx \num{2.1e3}$
个\omterm{def:g10:orders-of-magnitude:lightyear}{天文单位}。点题：把宇宙缩小，直到太阳只有一粒沙子那么大，
可观测宇宙仍有真实日地距离的两千倍之大——
没有任何比例模型能把宇宙缩小到人类尺度；
只有 10 的幂次才能承载它。
\end{solution}
